% Second pass (Claude, 2026-11-10): first pass checked against the archive images; changes in LM_C3-568-586_changes.md.
\scanpage{568}
% source-position: top, lines 1--2
immediate from (77). \deleted{It therefore remains to}\inserted{We thus have}
\deleted{\unsure{d}}

% source-position: upper, display (79)
\[
(79)\qquad
\diagram{\begin{tikzcd}
S_0\Gamma_{\rho,\sigma} \arrow[r,hook] & S\unsure{\mathfrak{G}}_{\rho,\sigma}^{\mathrm{eff}}
\end{tikzcd}}
\]
% source-position: upper, lines 3--5
it remains to see that this map is surjective ---
we have done what was needed for. Let us first note
that the composite
\[
\diagram{\begin{tikzcd}
S_0\Gamma_{\rho,\sigma} \arrow[r,hook] & S\unsure{\mathfrak{G}}_{\rho,\sigma}^{\mathrm{eff}} \arrow[r,"\text{epimorphism}"] &[30pt] \Sigma_{\rho,\sigma}^{\mathrm{eff}}\bigl( \arrow[r,"\sim"] & Z_{\rho,\sigma}\bigr)
\end{tikzcd}}
\]
\deleted{\(\diagram{\begin{tikzcd}(u,a,b)\arrow[r,mapsto]&u\end{tikzcd}}\)}
% source-position: upper-middle, lines 6--7
is none other, by (77) and (71) than
\deleted{\(\diagram{\begin{tikzcd}
(\unsure{\alpha},\unsure{\beta},\unsure{\mu})\arrow[r,mapsto]&
(\unsure{[\alpha,\rho]},\unsure{[\beta,\sigma]})
\end{tikzcd}}\)}
\[
(80)\qquad
\diagram{\begin{tikzcd}
x=(u,a,b)\arrow[r,mapsto]& (\xi,\eta,\mu,\varepsilon',\varepsilon)
\end{tikzcd}}
\quad
\left\{\begin{array}{ll}
\xi=[u,\rho]&\text{thus }u(\rho)=\xi\rho\\
\eta=[u,\sigma]&\text{thus }u(\sigma)=\eta\sigma\\
\mu=\chi_{S\Gamma}(x)&\\
\varepsilon'=\varepsilon'(x)&\\
\varepsilon=\varepsilon(x)&
\end{array}\right.
\]
% source-position: middle, lines 8--10
thus its \deleted{image} composite with
\(\diagram{\begin{tikzcd}\Sigma_{\rho,\sigma}^{\mathrm{eff}}\arrow[r,"\sim"]&Z_{\rho,\sigma}\end{tikzcd}}\)
is none other than the canonical projection
\[
\diagram{\begin{tikzcd}
S_0\Gamma_{\rho,\sigma}\arrow[r]&Z_{\rho,\sigma}\\
(u,a,b)\arrow[r,mapsto]&u
\end{tikzcd}}
\]
which is surjective, with kernel
\(SN_\rho^*\times SN_\sigma^*=SZ_\rho\times SZ_\sigma\).

% source-position: lower-middle, lines 11--12
Let us consider on the other hand \deleted{\unsure{the operation}}
\deleted{It remains to see that}
\[
(81)\quad
\left\{\begin{array}{l}
\text{operation of }SN_\rho\times SN_\sigma\\
\text{on }\deleted{S}\unsure{\mathfrak{G}}_{\rho,\sigma}
\subset\mathcal{G}\times\mathcal{G}\times\widehat{\mathbb Z}^{*}\\
\hfill\Vert\\
\{(\alpha,\beta,\mu)\mid[\alpha,\rho]=[\beta,\sigma]\varepsilon_1^{\mu-1}\}
\end{array}\right.
\quad
\begin{array}{l}
(a_0,b_0)\in S\deleted{\unsure{N}}Z_\rho^{\unsure{*}}\times S\deleted{\unsure{N}}Z_\sigma^{\unsure{*}}\\
\inserted{(=SZ_\rho\times SZ_\sigma)}\\
\text{operates by}\\
\diagram{\begin{tikzcd}(\alpha,\beta,\mu)\arrow[r,mapsto]&(\alpha a_0,\beta b_0,\mu)\end{tikzcd}}
\end{array},
\]
% source-position: lower, lines 13--17
It is clear that it operates \deleted{faithfully} freely, and
that the quotient of \(\unsure{\mathfrak{G}}_{\rho,\sigma}\) by this action
is identified with \(\Sigma_{\rho,\sigma}\), thus we have the result
idem for \(\unsure{\mathfrak{G}}_{\rho,\sigma}^{\unsure{*}}\) and
\(\Sigma_{\rho,\sigma}^{\unsure{*}}\) \deleted{\illegible{1}}. Now, for
% source-position: bottom, overwritten display (82), lines 18--20
\deleted{(82) for \(\varepsilon,\varepsilon'\)
\inserted{\(\varepsilon=(-1)^\nu,\quad\varepsilon'=(-1)^{\nu'}\)} fixed,
the action (81) of \(SZ'_\rho\times SZ'_\sigma\) corresponds
in \(S_0\Gamma_{\rho,\sigma}\) to the \deleted{translation}\inserted{operation}
\unsure{\((u,a,b)\)}
\(\diagram{\begin{tikzcd}
(u,a,b)\arrow[r,mapsto]&
\bigl(u,\unsure{\varepsilon_\rho^{\tau'}(a_0^{-1})}a,
\unsure{\varepsilon_\sigma^{\tau'}(b_0^{-1})}b\bigr)
\end{tikzcd}}\)}
\mspage{567}

% Printed archive footer outside the manuscript frame (not manuscript):
% 567
% https://grothendieck.umontpellier.fr
% Université de Montpellier, UMR CNRS 5149, Case courrier 051, Place Eugène Bataillon
% 34095 Montpellier cedex 5 - France

\scanpage{569}
% source-position: top edge, marginal calculation
\margin{\deleted{\unsure{i.e.}\ \(\unsure{\varepsilon_\rho^{\tau'}}\,
\unsure{(a)}=\tau'^{\nu'}\)}}
\mspage{588}
% source-position: top line 1, cancelled
\deleted{\unsure{it is thus} (\illegible{2} \unsure{the set}}
% source-position: upper display (82)
\[
(82)\qquad \varepsilon=(-1)^\nu,\quad\varepsilon'=(-1)^{\nu'}
\]
fixed, the preceding operation of \((a_0,b_0)\) corresponds,
under the map (76), to the operation
\[
(83)\quad
\diagram{\begin{tikzcd}
(u,a,b)\arrow[r]&\bigl(u,\tau^{\nu'}(a_0^{-1})a,\tau^\nu(b_0^{-1})b\bigr)
\end{tikzcd}},
\]
% source-position: upper-middle, lines 5--8
i.e. to left translation by
\((\tau^{\nu'},\tau^\nu)\underbrace{(a_0^{-1},b_0^{-1})}_{\inserted{(a_0,b_0)^{-1}}}\).
Whence surjectivity, and the desired isomorphism
\[
(84)\qquad
\diagram{\begin{tikzcd}
S_0\Gamma_{\rho,\sigma}\arrow[r,"\sim"]&S\unsure{\mathfrak{G}}_{\rho,\sigma}^{\mathrm{eff}}
\end{tikzcd}}
\quad\text{(cf. 76)}
\]
% source-position: middle, lines 9--12
The left translation by \(\ell\in L_0\)
\inserted{in \(S_0\Gamma_{\rho,\sigma}\)} corresponds,
under the preceding bijection, to the \deleted{translation}\inserted{operation}
\[
(85)\qquad
\diagram{\begin{tikzcd}
(\alpha,\beta,\mu)\arrow[r,mapsto]&(\ell\alpha,\ell\beta,\mu)
\end{tikzcd}}
\]
which it was clear a priori stabilizes
\(S\unsure{\mathfrak{G}}_{\rho,\sigma}^{\mathrm{eff}}\).
Passing to the quotient, we thus find
\[
(86)\qquad
\diagram{\begin{tikzcd}
\Gamma_{\rho,\sigma}\arrow[r,"\sim"]&L_0\backslash S\unsure{\mathfrak{G}}_{\rho,\sigma}^{\mathrm{eff}}
\end{tikzcd}}.
\]

% source-position: lower, lines 13--14 and display (87)
Finally, let us consider
\[
(87)\quad
\left\{\begin{array}{l}
\diagram{\begin{tikzcd}[column sep=large]
S_0M_{\rho,\sigma}\arrow[r]&S\unsure{\mathfrak{G}}_{\rho,\sigma}^{\mathrm{eff}}\times\mathcal{G}^+
=\underbrace{\{\overset{\inserted{(\in\mathcal{G}\times\mathcal{G}\times\widehat{\mathbb Z}^{*}\times\mathcal{G}^{+}}}{(\alpha,\beta,\mu,f)}\mid}
_{\unsure{(}[\alpha,\rho]=[\beta,\sigma]\varepsilon_1^{\mu-1}\unsure{)}}\\
x=(u,a,b,U)\arrow[r,mapsto]&
\deleted{((\unsure{ua^{-1}\tau'},\unsure{ub^{-1}\tau'}),\unsure{uU^{-1}})}
\inserted{((\alpha,\beta,\mu),f)}
\end{tikzcd}}
\\
\hfill
\left\{\begin{array}{l}
\alpha=ua^{-1}\tau'^{\deleted{\unsure{\nu'}}\inserted{\unsure{\nu'_\rho(a)}}}\\
\beta=ub^{-1}\tau'^{\unsure{\nu_\sigma(b)}}\\
\mu=\chi_{SM}(x)\quad(=\chi_Z(u)=\chi_\rho(a)=\chi_\sigma(b)=\chi_G(U))\\
f=uU^{-1}
\end{array}\right.
\end{array}\right.
\]
% source-position: bottom line 15
This map is bijective (trivial by (84)).
\mspage{568}

% Printed archive footer outside the manuscript frame (not manuscript):
% 568
% https://grothendieck.umontpellier.fr
% Université de Montpellier, UMR CNRS 5149, Case courrier 051, Place Eugène Bataillon
% 34095 Montpellier cedex 5 - France

\scanpage{570}
\mspage{588}
% source-position: top heading
\textcircled{\(7^\circ\)}\quad\uline{Functorialities}

Let us consider
\[
(88)\qquad
\diagram{\begin{tikzcd}\mathcal{G}\arrow[r,"\varphi"]&\mathcal{G}'\end{tikzcd}}
\qquad\begin{array}{l}\text{surjective homomorphism of}\\\text{profinite groups}\end{array}
\]
% source-position: upper lines 3--5
let \deleted{\(\unsure{\rho,\sigma,\tau'}\)}
\(\rho_{\mathcal{G}'},\sigma_{\mathcal{G}'},\tau'_{\mathcal{G}'},\ldots\)
be the images of \(\rho=\rho_{\mathcal{G}},\sigma=\sigma_{\mathcal{G}},
\tau'=\tau'_{\mathcal{G}},\ldots\) in \(\mathcal{G}'\).

% source-position: upper-middle line 6
\uline{Proposition}\quad Equivalent conditions on \(\varphi\) :

a) The map induced by \(\varphi\)
\[
(89)\qquad
\diagram{\begin{tikzcd}
\Sigma_{\rho,\sigma}^{\mathcal{G}}\arrow[r]&\Sigma_{\rho,\sigma}^{\mathcal{G}'}
\end{tikzcd}}
\]
\deleted{\unsure{transforms}}\inserted{maps}
\(\Sigma_{\rho,\sigma}^{\mathcal{G}\,\mathrm{eff}}\)
into \(\Sigma_{\rho,\sigma}^{\mathcal{G}'\,\mathrm{eff}}\).

% source-position: middle lines 8--11
b) For every \(u\in\operatorname{Aut}\mathcal{G}^+\),
\unsure{which|that} normalizes \(L_0\),
and \unsure{such|such} that \(u(\deleted{\rho}L_\rho)\underset{\mathcal{G}^+}{\sim}L_\rho\),
\(u(L_\sigma)\underset{\mathcal{G}^+}{\sim}L_\sigma\),
there exists \(u'\in\operatorname{Aut}(\mathcal{G}'^+)\) which makes commutative
\[
(90)\qquad
\diagram{\begin{tikzcd}
\mathcal{G}^+\arrow[r,"\varphi"]\arrow[d,"u"']&\mathcal{G}'^+\arrow[d,"u'","{\unsure{-}}"']\\
\mathcal{G}^+\arrow[r,"\varphi"]&\mathcal{G}'^+
\end{tikzcd}},
\]
in other words, \(u\) \deleted{\unsure{maps}} \deleted{\unsure{stabilizes}}
\inserted{\unsure{respects}} \(\operatorname{Ker}\varphi\).

% source-position: middle-lower lines 12--15
b') idem, assuming only, in place of
\(u\) normalizing \(L_0\), that
\(u(L_0)\underset{\mathcal{G}^+}{\sim}L_0\)
(plus the other conditions)

\unsure{c|b}) \(\exists\) group homomorphism
\[
(91)\qquad
\varphi_Z:\quad
\diagram{\begin{tikzcd}
Z_{\rho,\sigma}^{\mathcal{G}}\arrow[r,"{\deleted{\unsure{\varphi_Z}}}"]&Z_{\rho,\sigma}^{\mathcal{G}'}
\end{tikzcd}}
\]
\deleted{\unsure{with} \(\unsure{u'}\)}
% source-position: lower lines 16--19
compatible with \(\chi,\varepsilon_Z,\varepsilon'_Z\) and such that
\(\forall u\in Z_{\rho,\sigma}^{\mathcal{G}}\), if
\(u'\) is its image we have commutativity in (90).
(\uline{N.B.} \(\varphi_Z\) \unsure{will be} \deleted{\illegible{1}}
unique \unsure{by|for} \unsure{\(\zeta\)|?}, and is compatible with \(i_{L_0,Z}\))
\margin{thus it induces
\(\diagram{\begin{tikzcd}
\gamma_{\rho,\sigma}^{\mathcal{G}}\arrow[r]&\gamma_{\rho,\sigma}^{\mathcal{G}'}
\end{tikzcd}}\)}

% source-position: bottom lines 20--23
c') \(\exists\) group homomorphism
\[
(92)\qquad
\varphi_{\deleted{Z}\inserted{G}}:\quad
\diagram{\begin{tikzcd}
G_{\rho,\sigma}^{\mathcal{G}}\arrow[r]&G_{\rho,\sigma}^{\mathcal{G}'}
\end{tikzcd}}
\]
compatible with \(\chi_G,\varepsilon_G,\varepsilon'_G\) and satisfying the
compatibility condition (90). (N.B. this \(\varphi_G\)
\mspage{569}

% Printed archive footer outside the manuscript frame (not manuscript):
% 569
% https://grothendieck.umontpellier.fr
% Université de Montpellier, UMR CNRS 5149, Case courrier 051, Place Eugène Bataillon
% 34095 Montpellier cedex 5 - France

\scanpage{571}
\mspage{589}
% source-position: top lines 1--4
is then unique, and compatible with \(i_{\mathcal{G},G}\)
\inserted{\(i_{Z,G}\)}\unsure{;}
\deleted{\unsure{transforms}}\inserted{maps}
\(\deleted{\unsure{N}}\inserted{Z}_\rho^{\mathcal{G}}\) into
\(\deleted{\unsure{N}}\inserted{Z}_\rho^{\mathcal{G}'}\),
\(\deleted{\unsure{N}}\inserted{Z}_\sigma^{\mathcal{G}}\) into
\(\deleted{\unsure{N}}\inserted{Z}_\sigma^{\mathcal{G}'}\),
\unsure{which} extends to
\(\varphi_\rho:\diagram{\begin{tikzcd}N_\rho^{\mathcal{G}}\arrow[r]&N_\rho^{\mathcal{G}'}\end{tikzcd}}\)
inducing
\(\diagram{\begin{tikzcd}N_\rho^{*\mathcal{G}}\arrow[r]&N_\rho^{*\mathcal{G}'}\end{tikzcd}}\),
and extends to
\(\varphi_\sigma:\diagram{\begin{tikzcd}N_\sigma^{\mathcal{G}}\arrow[r]&N_\sigma^{\mathcal{G}'}\end{tikzcd}}\)
inducing
\(\diagram{\begin{tikzcd}N_\sigma^{*\mathcal{G}}\arrow[r]&N_\sigma^{*\mathcal{G}'}\end{tikzcd}}\)\ldots)

% source-position: upper lines 5--7
\deleted{\unsure{it} \illegible{4}}
\deleted{d) \(\exists\) homomorphism of groups}

We shall then say that \(\varphi\) is ``extendable''.

% source-position: upper-middle lines 8--10
\uline{Corollary}\quad We can then extend \(\varphi\) to
\[
(93)\qquad\varphi_M:
\diagram{\begin{tikzcd}
S_0G_{\rho,\sigma}^{\mathcal{G}}\arrow[r]&S_0G_{\rho,\sigma}^{\mathcal{G}'}
\end{tikzcd}},
\]
compatible with \(\Theta_Z,\Theta_\rho,\Theta_\sigma,\Theta_G\)
(taking into account \(\varphi_Z,\varphi_\rho,\varphi_\sigma,\varphi_G\)),
and this in a way ``compatible with \unsure{everything}'', in particular
with diagrams (45) (62).

% source-position: middle-lower lines 11--14
\uline{Example}\quad Suppose
\(\Sigma_{\rho,\sigma}^{\mathcal{G}'}=\Sigma_{\rho,\sigma}^{\mathcal{G}'\,\mathrm{eff}}\),
then the conditions of the proposition are \unsure{automatically}
satisfied (whatever the surjective homomorphism
\(\diagram{\begin{tikzcd}\mathcal{G}\arrow[r]&\mathcal{G}'\end{tikzcd}}\)).

% source-position: lower lines 15--19
\uline{Remark}\quad It may be that we are given a subgroup
\(Z_{\rho,\sigma}^{\unsure{\natural|\sharp}}\subset Z_{\rho,\sigma}^{\mathcal{G}}\),
and that we consider \deleted{\unsure{a)}} condition a) of the proposition
only for \(u\in Z_{\rho,\sigma}^{\unsure{\natural|\sharp}}\).
But with the help of \(Z_{\rho,\sigma}^{\unsure{\natural|\sharp}}\), we construct
\margin{We must suppose \unsure{in addition} that
\(\tau'_Z\in Z_{\rho,\sigma}^{*}\)}
\[
(94)\quad
G_{\rho,\sigma}^{\unsure{\natural|\sharp}}=
Z_{\rho,\sigma}^{\unsure{\natural|\sharp}}\cdot\mathcal{G}^+,
\quad Z_\rho^{\unsure{\natural|\sharp}}=Z_\rho\cap G_{\rho,\sigma}^{\unsure{\natural|\sharp}},
\quad Z_\sigma^{\unsure{\natural|\sharp}}=Z_\sigma\cap G_{\rho,\sigma}^{\unsure{\natural|\sharp}}
\]
% source-position: bottom lines 20--22
whence \(N_\rho^{\unsure{\natural|\sharp}}\deleted{_{\unsure{\rho,\sigma}}}\)
and \(N_\rho^{\unsure{\natural|\sharp}*}\deleted{_{\unsure{\rho,\sigma}}}\),
\(N_\sigma^{\unsure{\natural|\sharp}}\) and
\(N_\sigma^{\unsure{\natural|\sharp}*}\),
\(\gamma_{\rho,\sigma}^{\unsure{\natural|\sharp}}\), then
\(\unsure{S}_0M_{\rho,\sigma}^{\unsure{\natural|\sharp}}\) and
\(M_{\rho,\sigma}^{\unsure{\natural|\sharp}}\).
\deleted{We have a variant of \unsure{the}}
This being so, we have an obvious variant of the
\mspage{570}

% Printed archive footer outside the manuscript frame (not manuscript):
% 570
% https://grothendieck.umontpellier.fr
% Université de Montpellier, UMR CNRS 5149, Case courrier 051, Place Eugène Bataillon
% 34095 Montpellier cedex 5 - France

\scanpage{572}
% source-position: top lines 1--2
proposition, giving conditions for \unsure{extension|extension} of \(\varphi\) to
\(\diagram{\begin{tikzcd}
S_0M_{\rho,\sigma}^{*}\arrow[r]&S_0M_{\rho,\sigma}^{\mathcal{G}'}
\end{tikzcd}}\).

% source-position: upper abandoned block, lines 3--7
\deleted{\uline{Universal case}
\[
\mathcal{G}=\operatorname{Gl}(2,\mathbb Z)^\wedge,\quad
\rho=\begin{pmatrix}0&1\\-1&1\end{pmatrix},\quad
\sigma=\begin{pmatrix}0&-1\\1&0\end{pmatrix}
\]
\[
\tau'=\begin{pmatrix}0&1\\1&0\end{pmatrix},\quad
\unsure{\{}\ \tau=\begin{pmatrix}-1&0\\0&+1\end{pmatrix};\quad
\varepsilon_0=\begin{pmatrix}1&-1\\0&1\end{pmatrix},
\]
\deleted{We have}
\[
\varepsilon_1=\begin{pmatrix}\unsure{1}&0\\1&1\end{pmatrix}
\deleted{^{\unsure{\varepsilon}}}.
\]
Since \(\rho\)}
\mspage{571}

% Printed archive footer outside the manuscript frame (not manuscript):
% 571
% https://grothendieck.umontpellier.fr
% Université de Montpellier, UMR CNRS 5149, Case courrier 051, Place Eugène Bataillon
% 34095 Montpellier cedex 5 - France

\scanpage{573}
\mspage{590}
% source: upper portion, heading, lines 1--5, (1)--(2).

$8^\circ)$ \uline{The universal case} $\mathcal{G}=\mathrm{Gl}(2,\mathbb{Z})^{\wedge}$.

\unsure{.} Let us consider the exact sequence

\diagram{\begin{tikzcd}
\text{(1)}\quad 1 \arrow[r] & \{\pm1\} \arrow[r] & \mathcal{N}_{1,1}^{\sim} \arrow[r] \arrow[d,equal] & \mathcal{M}_{0,3}^{\sim} \arrow[r] & 1 \\
 & & \mathcal{M} & &
\end{tikzcd}}

\deleted{the group $\{\pm1\}$}
We have
\[
(2)\qquad \{\pm1\}=\unsure{\mathrm{Centr}|\mathrm{Centre}}(\mathcal{M})=\unsure{\mathrm{Centr}|\mathrm{Centre}}(\widehat{\mathcal{G}}^{+}),
\]

% source: oblique left margin beside (2)--(3), four identifications.
\margin{\unsure{we have taken|we have set}
\[
\begin{aligned}
\mathcal{G}&=\unsure{\mathcal{T}_{1,1}|\mathcal{G}_{1,1}}\simeq\mathrm{Gl}(2,\mathbb{Z}),\\
\widehat{\mathcal{G}}&=\widehat{\mathcal{T}}_{1,1}\simeq\mathrm{Gl}(2,\mathbb{Z})^{\wedge},\\
\unsure{\widehat{\mathcal{G}}^{+}|\mathcal{G}^{+}}&=\mathcal{T}_{1,1}^{+}=\mathrm{Sl}(2,\mathbb{Z}),\\
\widehat{\mathcal{G}}^{+}&=\mathrm{Sl}(2,\mathbb{Z})^{\wedge}.
\end{aligned}
\]}

% source: middle portion, lines 6--12, (3)--(4).
whence by (1) a map

\diagram{\begin{tikzcd}
\text{(3)}\quad \mathcal{M}_{0,3}^{\sim} \arrow[r] & \mathrm{Aut}(\mathcal{M}) \arrow[r,"{\inserted{\substack{\text{restriction to}\\ \mathcal{G}^{+\wedge}=\mathrm{Sl}(2,\mathbb{Z})^{\wedge}}}}"] &[110pt] \mathrm{Aut}(\widehat{\mathcal{G}}^{+\deleted{\unsure{\wedge}}})
\end{tikzcd}}

\deleted{in the case}
Besides, \unsure{every} \unsure{automorphism} of $\widehat{\mathcal{G}}^{+}$ \unsure{preserves}
its center, hence \unsure{passes} to the quotient, whence

\diagram{\begin{tikzcd}
\text{(4)}\quad \mathrm{Aut}(\widehat{\mathcal{G}}^{+}) \arrow[r] & \mathrm{Aut}(\widehat{\mathcal{G}}^{+}/\!\pm1\simeq\deleted{\unsure{\mathcal{M}}}\mathcal{G}_{0,3}^{+\wedge})
\end{tikzcd}}

% source: lower-middle portion, lines 13--16, triangle (5).
and \unsure{the composite|the composite} of (3) with (4) is the canonical
homomorphism:

\diagram{\begin{tikzcd}
\text{(5)}\quad \mathcal{M}_{0,3}^{\sim} \arrow[rr] \arrow[dr] & & \mathrm{Aut}(\widehat{\mathcal{G}}^{+}) \arrow[dl] \\
& \mathrm{Aut}(\widehat{\mathcal{G}}_{0,3}^{+}=\widehat{\mathcal{G}}^{+}/\{\pm1\}) &
\end{tikzcd}}

This shows that (3) is injective, since
$\mathcal{M}_{0,3}^{\sim}\longrightarrow\mathrm{Aut}(\widehat{\mathcal{G}}_{0,3}^{+})$
is. \deleted{The map}

% source: bottom portion, lines 17--20, remark continuing on scan574.
\uline{Remark} The map
$\mathrm{Aut}(\widehat{\mathcal{G}}^{+})\longrightarrow\mathrm{Aut}(\mathcal{G}_{0,3}^{+\wedge})$
is not injective on the other hand, \deleted{\unsure{the}} the
elements of the kernel correspond to the
\unsure{homomorphisms}
$\widehat{\mathcal{G}}^{+}\longrightarrow\unsure{\mathrm{Centr}|\mathrm{Centre}}(\widehat{\mathcal{G}}^{+})=\{\pm1\}$,

\mspage{572}

% Printed archive footer outside the manuscript frame (not manuscript):
% 572
% https://grothendieck.umontpellier.fr
% Université de Montpellier, UMR CNRS 5149, Case courrier 051, Place Eugène Bataillon
% 34095 Montpellier cedex 5 - France

\scanpage{574}
% source: upper portion, lines 1--5, (6)--(7), including right marginal map.

or again to the homomorphisms \unsure{from}

\diagram{\begin{tikzcd}
\text{(6)}\quad \mathcal{G}_{\mathrm{ab}}^{+\wedge} \arrow[r] & \{\pm1\} & \text{or again}\quad (\mathcal{G}_{\mathrm{ab}}^{+})_{2} \arrow[r] & \{\pm1\}
\end{tikzcd}}

Now we see that
\[
\mathcal{G}_{\mathrm{ab}}^{+\wedge}
=\left\{\rho,\sigma\ \middle|\ \begin{gathered}
\rho^{6}=\sigma^{4}=[\rho,\sigma]=1\\
\rho^{3}=\sigma^{2}
\end{gathered}\right\}=\mathbb{Z}/12\mathbb{Z}
\]

\margin{\diagram{\begin{tikzcd}
\rho \arrow[r,mapsto] & 2\bmod12\\
\sigma \arrow[r,mapsto] & 3\bmod12
\end{tikzcd}}}

hence

\diagram{\begin{tikzcd}
& (\mathcal{G}_{\mathrm{ab}}^{+})_{2} \arrow[r,"\sim"] & \mathbb{Z}/2\mathbb{Z}\\
\text{(7)}\quad\left\{\vphantom{\begin{matrix}1\\1\\1\end{matrix}}\right. & \rho \arrow[r,mapsto] & 0\\
& \sigma \arrow[r,mapsto] & 1
\end{tikzcd}}

% source: middle portion, lines 6--9, composite diagram (8).
\deleted{hence to}
\inserted{This isomorphism corresponds to the character}
$\mathcal{G}^{+\wedge}\longrightarrow\{\pm1\}$ which is also obtained as the
composite

\diagram{\begin{tikzcd}[column sep=large]
\text{(8)}\quad \mathcal{G}^{+\wedge} \arrow[r] \arrow[rr,bend left=27,shorten >=50pt,"{\unsure{\mathrm{sg}|\mathrm{sgn}}}"{yshift=4pt}] &
\underset{\substack{\text{symmetric group}\\\text{on three letters}}}{\mathcal{G}_{3}}
=\mathcal{G}^{+\wedge}/\underset{\substack{\text{inverse}\\\text{image}\\\text{of }\Pi_{0,3}^{\wedge}\\\text{in }\mathcal{G}^{+\wedge}}}{\widehat{\Pi}}
\arrow[r,"{\substack{\text{sign}\\\text{\unsure{of a} permutation}}}"] &[50pt] \{\pm1\}
\end{tikzcd}}

% source: middle-lower portion, lines 10--13, (9)--(10).
Thus the only element \inserted{$\Theta$} \deleted{\unsure{different}} from $1$
of the kernel of \deleted{(3) is given by} \inserted{(4) is given by}

\diagram{\begin{tikzcd}
\text{(9)}\quad \deleted{\mathrm{Aut}\mathcal{G}}\ \inserted{\Theta:}\ x \arrow[r,mapsto] & x\cdot\mathrm{sg}(x)
\end{tikzcd}}

We shall note, \deleted{that \unsure{$\Theta$ and}} on setting
\[
(10)\qquad L_{0}^{\mathcal{G}}=\varepsilon_{0}^{\widehat{\mathbb{Z}}},\qquad
L_{0}^{\Pi}=L_{0}^{\mathcal{G}}\cap\Pi=l_{0}^{\widehat{\mathbb{Z}}},
\]
% source: lower portion, heavily rewritten lines 14--17, (11), ending unfinished.
that the subgroup $\Theta(L_{0}^{\mathcal{G}})$\inserted{$\subset\mathcal{G}^{+\wedge}$}
is \uline{not} $\mathcal{G}$-conjugate to $L_{0}$,
\deleted{\unsure{or}} \deleted{and $\Theta(L_{0}^{\Pi})$ $\mathcal{G}$-conjugate}
\inserted{for} \inserted{it is generated by}
\deleted{\unsure{$\Theta\varepsilon$}}\deleted{and a fortiori not.} :
\[
(11)\qquad \Theta(\varepsilon_{0})=-\varepsilon_{0}
\qquad\left(\begin{gathered}
\text{for }\varepsilon_{0}=\sigma\rho\\
\mathrm{sg}(\varepsilon_{0})=-1
\end{gathered}\right)
\]
and we check that $\varepsilon_{0}$ is not conjugate to
$(-\varepsilon_{0})^{\mu}=\deleted{\unsure{\varepsilon_{0}}}-\varepsilon_{0}^{\mu}$
for \unsure{any} $\mu\in\widehat{\mathbb{Z}}^{*}$ (as we see

\mspage{573}

% Printed archive footer outside the manuscript frame (not manuscript):
% 573
% https://grothendieck.umontpellier.fr
% Université de Montpellier, UMR CNRS 5149, Case courrier 051, Place Eugène Bataillon
% 34095 Montpellier cedex 5 - France

\scanpage{575}
\mspage{591}
% source: upper portion, pale continuation, lines 1--9 and (12).

at once, on \deleted{passing: $\mathcal{G}/\Pi_{\unsure{a}}$}\inserted{passing:}
$\mathcal{G}_{0,3}^{\wedge}/\Pi_{0,3}^{\wedge}{}_{\mathrm{ab}}$
and \unsure{noting} that $l_{0}=\varepsilon_{0}^{2}$ would be
conjugate to: $(-\varepsilon_{0}^{\mu})^{2}=l_{0}^{\mu}$,
\unsure{that we must have|which must have from it} $\mu=1$, now
$\varepsilon_{0}$ and $-\varepsilon_{0}$ are not conjugate
even in $\mathrm{Sl}(2,\widehat{\mathbb{Z}})$\unsure{'|!} ... ). Now we shall
be interested \unsure{precisely} in the automorphisms \unsure{of|$\varphi$}
$\mathcal{G}^{+\wedge}$ which transform $L_{0}^{\mathcal{G}}$ into a
conjugate subgroup. We \unsure{shall note|note} that on the other hand
\[
(12).\qquad \Theta(l_{0})=l_{0}\qquad\text{hence}\qquad
\Theta(L_{0}^{\Pi})=L_{0}^{\Pi}
\]

% source: middle portion, proposition and conditions a)--b), lines 10--17.
\Needspace{12\baselineskip}
\uline{Proposition} The homomorphism (3)

\diagram{\begin{tikzcd}
\mathcal{M}_{0,3}^{\sim} \arrow[r] & \mathrm{Aut}(\widehat{\mathcal{G}}^{+})
\end{tikzcd}}

is injective, and its image consists of the
automorphisms\inserted{$u$} of $\widehat{\mathcal{G}}^{+}$ which satisfy the following
two conditions

a) $u$ stabilizes $\widehat{\Pi}=\mathcal{G}^{+\wedge}[2]
=\mathrm{Ker}(\mathcal{G}^{+\wedge}\longrightarrow\mathrm{Sl}(2,\mathbb{Z}/2\mathbb{Z}))$

b) $u$ transforms $L_{0}^{\mathcal{G}}$ into a subgroup of $\mathcal{G}^{+\wedge}$
$\mathcal{G}^{+\wedge}$-conjugate to $L_{0}$; i.e.
$\exists\mu\in\widehat{\mathbb{Z}}^{*}$ and $f_{0}\in\mathcal{G}^{+\wedge}$
such that we have
\[
(13)\qquad u(\varepsilon_{0})=\mathrm{int}(f_{0}^{-1})(\varepsilon_{0}^{\mu}).
\]

% source: lower portion, proof opening, lines 18--21.
\uline{Proof.} Injectivity has already been seen.
\deleted{\unsure{Let us prove|Let us recall}}\inserted{Let us show}
that conditions a)b) are necessary for
$u$ to \deleted{\unsure{belong}}\inserted{come}
from a $u_{0}\in\mathcal{M}_{0,3}^{\sim}$,
\deleted{For} a) this is trivial, since $u_{0}$ stabilizes
$\Pi_{0,3}^{\unsure{\wedge|\sim}}$, hence its

\mspage{574}

% Printed archive footer outside the manuscript frame (not manuscript):
% 574
% https://grothendieck.umontpellier.fr
% Université de Montpellier, UMR CNRS 5149, Case courrier 051, Place Eugène Bataillon
% 34095 Montpellier cedex 5 - France

\scanpage{576}
% source: upper portion, continuation of proof, lines 1--7 and (14).

inverse image in $\mathcal{G}^{+\wedge}$. For condition
b), we are reduced to the case where $u_{0}\in\mathcal{M}_{0,3}^{\sim}[0]$
(conjugating if necessary by an element of
$\mathcal{G}_{0,3}^{+\wedge}\simeq\mathcal{G}^{+\wedge}/\!\pm1$),
in this case we have by construction of $\mathcal{N}_{1,1}^{\sim}$ from
$\mathcal{M}_{0,3}^{\sim}$ (cf. \uline{\(\overline{\mathrm{III}}\)}, \unsure{pages|pages} 665...)
\[
(14)\qquad u_{0}(\varepsilon_{0})=\varepsilon_{0}^{\mu}
\]
where $\mu\in\widehat{\mathbb{Z}}^{*}$ is the multiplier, OK.

% source: middle portion, converse argument, lines 8--17.
Conversely, let $u$ be an automorphism of $\mathcal{G}^{+\wedge}$ which
satisfies a),b), consider the automorphism $\widetilde{u}$ of
$\mathcal{G}^{+\wedge}/\!\pm1\simeq\mathcal{G}_{0,3}^{+\wedge}$
which it induces. We thus see
that $\widetilde{u}(\Pi_{0,3}^{\wedge})=\Pi_{0,3}^{\wedge}$, and that
$\widetilde{u}(\varepsilon_{0})$ is $\mathcal{G}_{0,3}^{\wedge}$-conjugate
to an $\varepsilon_{0}^{\mu}$ ($\mu\in\unsure{\widehat{\mathbb{Z}}^{*}|\widehat{\mathbb{Z}}}$).
\deleted{Let} Thus by
definition, we have $\widetilde{u}\in\mathcal{M}_{0,3}^{\sim}$
(cf. \unsure{pages|pages} 390 ff).
Let $u'$ be the automorphism of $\mathcal{G}^{+\wedge}$ induced by $\widetilde{u}$
(via (3)) -- then $u'$ and $u$ induce the same
automorphism of $\mathcal{G}_{0,3}^{+\wedge}=\mathcal{G}^{+\wedge}/\{\pm1\}$,
thus we will have $u'=u\circ\Theta$, where
$\Theta$\inserted{$=u^{-1}u'$} is an automorphism which

% source: lower-middle portion, numbered properties and conclusion, lines 18--21.
$1^\circ)$ induces the identity on $\mathcal{G}^{+\wedge}/\{\pm1\}$, and

$2^\circ)$ takes $L_{0}^{\mathcal{G}}$ to a conjugate subgroup

-- we have seen that this implies $\Theta=1$, q.e.d.

% source: bottom portion, abandoned proposition and exact sequence (15).
\deleted{\uline{Proposition} Let $\widehat{\Pi}_{0}$ be the subgroup}
Let us now consider the exact \deleted{\unsure{lemma}}\inserted{sequence}

\diagram{\begin{tikzcd}
\text{(15)}\quad 1 \arrow[r] & \pm1 \arrow[r] & \widehat{\Pi} \arrow[r] \arrow[d,equal] & \Pi_{0,3}^{\wedge} \arrow[r] & 0\\
& & \mathcal{G}^{+\wedge}[2] & &
\end{tikzcd}}

\mspage{575}

% Printed archive footer outside the manuscript frame (not manuscript):
% 575
% https://grothendieck.umontpellier.fr
% Université de Montpellier, UMR CNRS 5149, Case courrier 051, Place Eugène Bataillon
% 34095 Montpellier cedex 5 - France

\scanpage{577}
\mspage{592}
% source: upper portion, lines 1--6, including (16)--(17).
and consider the splitting given by

\diagram{\begin{tikzcd}[column sep=large,row sep=large]
(16)\quad\left\{\vphantom{\begin{matrix}\Pi_{0,3}\\l_i\end{matrix}}\right.
& \Pi_{0,3} \arrow[r]
& \Pi=\mathcal{G}^{+}[2]=\operatorname{Ker}\bigl(\mathrm{Sl}(2,\mathbb Z)\longrightarrow\mathrm{Sl}(2,\mathbb Z/2\mathbb Z)\bigr) \\
& l_i \arrow[r,mapsto]
& \deleted{\unsure{\lambda_i|l_i}\overset{\unsure{\mathrm{dfn}}}{=}}\;\lambda_i\overset{\mathrm{dfn}}{=}-l_i
\end{tikzcd}}

There \deleted{\unsure{is}} is such a homomorphism, by virtue of
\[
(17)\qquad
\left\{\begin{aligned}
(-l_\infty)(-l_1)(-l_0)&=1\quad\text{in }\mathrm{Sl}(2,\mathbb Z)\quad\text{i.e.}\\
l_\infty l_1l_0&=-1
\end{aligned}\right.
\]
\margin{$)$}

and it is obvious that this is a splitting. We also likewise have a splitting for $\widehat\Pi$ over $\widehat\Pi_{0,3}$.

% source: middle portion, lines 7--12, definition (18).
We denote by $\Pi_0$ the image of $16)$, i.e.
\[
(18)\qquad
\left\{\begin{aligned}
\Pi_0&=\text{subgroup of }\Pi\inserted{=\mathcal{G}^{+}[2]}\text{ generated by}\\
&\lambda_0,\lambda_1,\lambda_\infty\qquad(\text{where }\lambda_i=-l_i)
\end{aligned}\right.
\]
and we thus have a subgroup $\widehat\Pi_0$ of $\widehat\Pi$.

% source: middle-lower portion, proposition a)--b), lines 13--18.
\uline{Proposition}\quad a) The subgroup $\Pi_0$ of $\Pi$ is \deleted{normal}\inserted{of index} $2$, and it is normal in $\mathcal{G}^{\unsure{\wedge|*}}$.

b) The subgroup $\widehat\Pi_0$ of $\widehat\Pi$ is \deleted{normal, \unsure{and}}\inserted{of index $2$}, and it is normal in $\mathcal M$ -- or, what amounts to the same, stable under the action of $\mathcal M/\pm1\simeq\mathcal M_{0,3}^{\sim}$.

% source: lower portion, proof and (19); right-hand additions at the same height.
\uline{Proof}\quad Since
\[
(19)\qquad
\left\{\begin{aligned}
\rho(\lambda_i)&=\lambda_{i+1}\\
\sigma(\lambda_0)&=\lambda_{1}^{\deleted{\unsure{-1}}},\quad
\sigma(\lambda_1)=\lambda_0^{\deleted{\unsure{-1}}}\\
\tau(\lambda_0)&=\lambda_0^{-1},\quad\tau(\lambda_1)=\lambda_1^{-1}
\end{aligned}\right.
\]
\margin{$\begin{aligned}
\sigma(\lambda_\infty)&=\sigma(\lambda_1\lambda_0)^{-1}=\lambda_1^{-1}\lambda_0^{-1}\\
\deleted{\unsure{\tau'}(\lambda_\infty)}&\deleted{=\lambda_\infty^{-1}}\\
\tau(\lambda_\infty)&=\tau(\lambda_1\lambda_0)^{-1}=\lambda_0\lambda_1
\end{aligned}$}

we see that $\Pi_0$ is \deleted{\unsure{stable}} normalized by $\rho,\sigma,\tau$, hence by $\mathcal G$. That it is of index $2$ is obvious. Thus $\widehat\Pi_0$ is of index $2$ in $\widehat\Pi$ and normalized by $\widehat{\mathcal G}$. To see that it is
\mspage{576}

% Printed archive footer outside the manuscript frame (not manuscript):
% 576
% https://grothendieck.umontpellier.fr
% Université de Montpellier, UMR CNRS 5149, Case courrier 051, Place Eugène Bataillon
% 34095 Montpellier cedex 5 - France

\scanpage{578}

% source-position: upper portion, continuation and display (20).
normalized by $\mathcal M$ in its \unsure{full} entirety, since
\[
\mathcal M=\mathcal M(0)\mathbin{.}\mathcal G^{+\wedge},
\]
we are reduced to verifying that it is normalized by
$\mathcal M(0)\simeq\mathcal M_{0,3}^{\sim}(0)$.
Now we \unsure{have}, for $u\in\mathcal M_{0,3}^{\sim}(0)$
\[
(20)\qquad
\left\{
\begin{aligned}
u(\rho)&=\operatorname{int}(\alpha)\rho\\
u(\sigma)&=\varepsilon_2\operatorname{int}(\beta)\sigma\\
u(\varepsilon_0)&=\varepsilon_0^\mu
\end{aligned}
\right.
\qquad
\left\{
\begin{aligned}
\varepsilon_2&=\varepsilon_2(\mu)\\
\alpha,\deleted{\unsure{\beta}}&\in\widehat{\Pi}\mathbin{.}\{1,\tau\}\\
\beta&\in\widehat{\Pi}
\end{aligned}
\right.
\]
% source-position: upper-middle portion, rewritten display (21).
\uline{whence}
\[
\deleted{u(\ell_0)=}\quad
\deleted{\unsure{(20)}}\quad
(21)\qquad
\left\{
\begin{aligned}
u(\ell_0)&=\ell_0^\mu\\
u(-\ell_0)&=(-\ell_0)^\mu
\quad\text{i.e.}\quad
u(\lambda_0)=\lambda_0^\mu
\quad(\text{since }(-1)^\mu=-1)
\end{aligned}
\right.
\]
\margin{$\unsure{)}$}

% source-position: middle portion, cancelled proof openings and lambda_1 calculation.
\deleted{It \unsure{remains} to \unsure{see}}
\deleted{\unsure{Indeed|Finally}}
thus $u(\lambda_0)\in\widehat{\Pi}_0$, it remains to prove that
$u(\lambda_1)\in\widehat{\Pi}_0$, \unsure{now}
\[
u(\lambda_1)
=u(\rho(\lambda_0))
=\underbrace{(u\rho)}_{\inserted{\deleted{\unsure{u}}(\operatorname{int}(\alpha)\rho)u}}(\lambda_0)
=\operatorname{int}(\alpha)(\rho)
\left(\underbrace{u(\lambda_0)}_{\inserted{\lambda_0^\mu}}\right)
=\deleted{\unsure{\operatorname{int}(\alpha)(\lambda_1^\mu)}}.
\]
\deleted{and we have \unsure{seen} that}
and since $\widehat{\Pi}_0$ is normalized by $\mathcal G^\wedge$
and in particular by $\deleted{\unsure{\rho}}\alpha\rho\alpha^{-1}$,
we indeed find $u(\lambda_1)\in\Pi_0$.

% source-position: lower portion, corollary (22).
\uline{Corollary} We have
\[
(22)\qquad
\left\{
\begin{aligned}
\Pi&\simeq\Pi_0\times\mu\\
\widehat{\Pi}&\simeq\widehat{\Pi}_0\times\mu
\end{aligned}
\right.
\]
where $\mu=\{\pm1\}=\operatorname{Centr}(\mathcal G^{+\wedge})$.

% source-position: bottom portion, cancelled and inserted proof of corollary.
\deleted{It suffices to \unsure{note} that the element $\unsure{-1}$
\unsure{of} $\mu$ \unsure{is not} in $\Pi_0$.}
\inserted{This is trivial, since we have a splitting of the central extension (15).}

\mspage{577}

% Printed archive footer outside the manuscript frame (not manuscript):
% 577
% https://grothendieck.umontpellier.fr
% Université de Montpellier, UMR CNRS 5149, Case courrier 051, Place Eugène Bataillon
% 34095 Montpellier cedex 5 - France

\scanpage{579}
\mspage{593}

% source-position: upper portion, introductory word and definitions (23)--(24).
\unsure{Let|Let}
\[
(23)\qquad
\left\{
\begin{aligned}
\mu_\tau&=\{1,\tau\}\subset\mathcal G
&D_\tau&=\mu_\tau\times\mu\subset\mathcal G\\
\mu_\sigma&=\{1,\unsure{\sigma|\tau'}\}\subset\mathcal G
&D_\sigma&=\mu_\sigma\times\mu\subset\mathcal G
\end{aligned}
\right.
\]
so that we have
\[
(24)\qquad
\begin{aligned}
D_\tau\cap\unsure{\mathcal G^+}
&=D_\tau\cap\Pi=\mu
&D_\tau\cap\Pi_0&=\{1\}\\
D_\sigma\cap\unsure{\mathcal G^+}
&=D_\sigma\cap\Pi=\mu
&D_\sigma\cap\Pi_0&=\{1\},
\end{aligned}
\]
\margin{\unsure{idem} on replacing $\mathcal G,\Pi$ by
$\mathcal G^\wedge,\widehat{\Pi}$}
% source-position: upper-middle portion, definitions (25)--(26).
We set
\[
(25)\qquad
\left\{
\begin{aligned}
\unsure{\Pi_\tau}&=\mu_\tau\mathbin{.}\Pi=D_\tau\mathbin{.}\Pi_0,
&\Pi_{0\tau}&=\mu_\tau\mathbin{.}\Pi_0\\
\Pi_\sigma&=\mu_\sigma\mathbin{.}\Pi=D_\sigma\mathbin{.}\Pi_0,
&\Pi_{0,\sigma}&=\mu_\sigma\mathbin{.}\Pi_0
\end{aligned}
\right.
\]
so that we have
\[
(26)\qquad
\left\{
\begin{aligned}
\Pi_\tau&=\Pi_{0\tau}\times\mu\\
\Pi_\sigma&=\Pi_{0,\sigma}\times\mu
\end{aligned}
\right.,
\]
and idem on putting in $\wedge$'s.

% source-position: middle-lower portion, normalization and (27).
\uline{Normalization of $\alpha$}, $\beta$\quad
The relations $20$, determine $\alpha,\beta$ only ``modulo sign''.
We remove this ambiguity by taking
\[
(27)\qquad
\alpha\in\Pi_{0,\tau},\qquad\beta\in\Pi_0.
\]

% source-position: lower portion, quotient presentation (28).
\uline{The group} $\mathcal G/\Pi_0$.
It is described by generators $\rho,\sigma,\tau^{\unsure{\prime}}$ and relations \unsure{by}
\[
(28)\qquad
\left\{
\begin{aligned}
&\rho^6=\sigma^4=\tau^2=1,\quad \rho^3=\sigma^2\\
&\tau(\rho)=\sigma(\rho^{-1}),\quad \tau(\sigma)=\sigma^{-1}\\
&\tau(\rho)=\rho
\end{aligned}
\right.
\]
% source-position: bottom portion, calculation (29) and revised concluding relations.
For we have
\[
(29)\qquad
\tau(\rho)=(-\ell_1^{-1})\rho
\]
thus the last relation $28$ reads
$\deleted{\tau(\rho)=\rho}\inserted{-\ell_1^{-1}=1}$
i.e. $\lambda_1=1$
(and implies
$\deleted{\unsure{-\ell_\infty=-\ell_0}}\inserted{\lambda_\infty=\lambda_0}=1$
on applying $\rho$).

\mspage{578}

% Printed archive footer outside the manuscript frame (not manuscript):
% 578
% https://grothendieck.umontpellier.fr
% Université de Montpellier, UMR CNRS 5149, Case courrier 051, Place Eugène Bataillon
% 34095 Montpellier cedex 5 - France

\scanpage{580}

% source-position: upper portion, presentations (30)--(31).
In terms of $\rho,\sigma,\tau'$ we find likewise
\[
(30)\qquad
\left\{
\begin{aligned}
&\rho^6=\sigma^4=\tau'^2=1,\quad \rho^3=\sigma^2\\
&\tau'(\rho)=\rho^{-1},\quad\tau'(\sigma)=\sigma^{-1}\\
&\sigma(\rho)=\rho^{-1}
\end{aligned}
\right.
\]
The subgroup $\mathcal G^+/\Pi_0$ has the presentation in $\rho,\sigma$
\[
(31)\qquad
\left\{
\begin{aligned}
&\rho^6=\sigma^4=1,\quad \rho^3=\sigma^2\\
&\sigma(\rho)=\rho^{-1}
\end{aligned}
\right\}
\]

% source-position: upper-middle portion, action question and extension diagram (32).
What is the operation of $\mathcal M_{0,3}^{\sim}\simeq\mathcal M/\pm1$
on $\mathcal G^+/\Pi_0$ ?
\deleted{\unsure{it is an} operation}
\inserted{We have a structure} of extension
\[
(32)\qquad
\diagram{
\begin{tikzcd}
1 \arrow[r] & \mu \arrow[r] \arrow[d,equal]
& \mathcal G^+/\Pi_0 \arrow[r]
& \mathcal G^+/\Pi \arrow[r] \arrow[d,equal]
& 1\\
& \pm1 && \mathcal G_3 &
\end{tikzcd}
}
\]
% source-position: middle portion, prose following (32).
and the \unsure{automorphisms} in question preserve this extension structure.
The \unsure{elements|elements} of
$\mathcal M_{0,3}^{!\sim}\simeq\mathcal M^!/\{\pm1\}$
operate by inducing the identity in the quotient $\mathcal G_3$,
and of \unsure{course} also in $\mu$ --- now such an automorphism of
the extension corresponds to a homomorphism
% source-position: lower-middle portion, character maps (33).
\[
(33)\qquad
\diagram{
\begin{tikzcd}
\mathcal G_3 \arrow[r] & \mu \arrow[d,equal]\\
& \pm1
\end{tikzcd}
}
\quad\text{i.e.}\quad
\diagram{
\begin{tikzcd}
\mathcal G_{3\,\mathrm{ab}} \arrow[r] & \mu
\end{tikzcd}
}
\]
% source-position: lower portion, sign character and conclusion (34).
now
\[
\diagram{
\begin{tikzcd}
\mathcal G_{3\,\mathrm{ab}}\arrow[r,"\uline{\mathrm{sg}}"]
&\{\pm1\}=\mu
\end{tikzcd}
}
\]
thus the group of automorphisms of $\mathcal G^+/\Pi_0$ inducing
the identity in $\mathcal G_3$ is $\simeq\mathbb Z/2\mathbb Z$,
a generator \inserted{$\Theta$} is given by
\[
(34)\qquad
\Theta(x)=\operatorname{sg}(x)\mathbin{.}x
\quad\text{thus}\quad
\left\{
\begin{aligned}
\Theta(\rho)&=\rho\\
\Theta(\sigma)&=-\sigma
\end{aligned}
\right.
\]

\mspage{579}

% Printed archive footer outside the manuscript frame (not manuscript):
% 579
% https://grothendieck.umontpellier.fr
% Université de Montpellier, UMR CNRS 5149, Case courrier 051, Place Eugène Bataillon
% 34095 Montpellier cedex 5 - France

\scanpage{581}
\mspage{594}

% source-position: upper portion, displays (35) and (35 bis).
Comparing with the formulas $(20)$, we find
\[
(35)\qquad
\deleted{\unsure{u}}
\left.u\right|_{\mathcal G^{+\wedge}/\widehat{\Pi}_0}
=\Theta^{\varepsilon_2(u)}:
\quad
\diagram{
\begin{tikzcd}
\dot\rho\arrow[r,mapsto]&\dot\rho\\
\dot\sigma\arrow[r,mapsto]&\varepsilon_2(u)\dot\sigma
\end{tikzcd}
}
\]
i.e.
\[
(35\text{ bis})\qquad
\begin{gathered}
[\dot u,\dot\sigma]\\[-0.5ex]
\unsure{\Vert}\\[-0.5ex]
\inserted{\dot u(\dot\sigma)\dot\sigma^{-1}}
\end{gathered}
=\varepsilon_2(u)^{\cdot}
\quad\text{in }\mathcal G/\Pi_0.
\]
% source-position: upper left slanted margin, commutator formula and warning.
\margin{This can also be written
$[\dot u,\dot x]=\operatorname{sg}(x)^{\unsure{\nu_2(u)}}$
in $\mathcal G^{+\wedge}/\Pi_0$
$(u\in\mathcal M^!,x\in\mathcal G^{+\wedge})$
(careful, this \unsure{goes wrong|does} \unsure{if we assume there}
\unsure{for} $u\in\mathcal M^!$ only $u\in\mathcal M\ldots$)}

% source-position: upper-middle portion, revised heading and definition (36).
\deleted{I would like}
Let us take up again the definition
\[
(36)\qquad
Z=Z_{\rho,\sigma}\stackrel{\mathrm{dfn}}=
\left\{
u\in\operatorname{Aut}(\mathcal G^{+\wedge})\ \left|\
\begin{aligned}
&u(\unsure{L_\rho})\
\inserted{\mathcal G^{+\wedge}\text{-}}\text{conjugate to }\
\deleted{\unsure{\ell_\rho}}L_\rho\\
&u(\unsure{L_\sigma})\
\inserted{\mathcal G^{+\wedge}}\text{-conjugate : }L_\sigma\\
&u(\unsure{L_0})=L_0
\end{aligned}
\right.
\right\}
\]
so that we shall have, for $u\in Z$
% source-position: middle portion, display (37) and parameter conditions.
\[
(37)\qquad
\left\{
\begin{aligned}
&u(\rho)\quad\mathcal G\text{-conjugate to }\rho^{\varepsilon'}\\
&u(\sigma)\quad\mathcal G\text{-conjugate to }\sigma^{\unsure{\varepsilon}}
\quad(=\varepsilon\sigma)\\
&u(\varepsilon_0)=\varepsilon_0^\mu
\end{aligned}
\right.
\]
with $\mu\in\widehat{\mathbb Z}^{*}$ and
$\varepsilon,\varepsilon'\in\{\pm1\}$ well determined,
(\unsure{for}
\diagram{
\begin{tikzcd}
\widehat{\mathbb Z}\arrow[r]&\mathcal G\\
n\arrow[r,mapsto]&\varepsilon_0^n
\end{tikzcd}
}
is injective, $\rho$ and $\rho^{-1}$ are not
$\mathcal G^{+\wedge}$-conjugate, nor $\sigma$ and $\sigma^{-1}$).
% source-position: lower-middle portion, proof of stabilization and l_0, l_1 calculations.
I say that such a $u$ comes from $\mathcal M_{0,3}^{\sim}$ ---
it remains to check, by virtue of the proposition..., that $u$ stabilizes
$\widehat{\Pi}$, which is generated by $\ell_0,\ell_1,\ell_\infty$.
Now we shall have
\[
u(\ell_0)=\ell_0^\mu\in\widehat{\Pi}
\]
and
\[
u(\ell_1)=
\underbrace{u\rho}_{\inserted{
\begin{gathered}
\deleted{\unsure{\operatorname{int}(f)\rho u}}\\
\deleted{\unsure{\operatorname{int}(f)\rho^{\varepsilon'}u}}\\
\operatorname{int}(f)(\rho^{\varepsilon'})u\\
f\in\mathcal G^{+\wedge}
\end{gathered}
}}(\ell_0)
=\operatorname{int}(f)(\rho^{\varepsilon'})(u(\ell_0))
=\operatorname{int}(f)(\rho)^{\varepsilon'}(\ell_0^\mu)
\]
now $\mathcal G^\wedge$ and a fortiori
$\operatorname{int}(f)(\rho)^{\varepsilon'}$ normalizes
$\deleted{\unsure{\mathcal G}}\widehat{\Pi}$,
thus \deleted{and} $u(\ell_1)\in\widehat{\Pi}$
% source-position: lower portion, l_infinity calculation and final cancelled words.
and since $u(-1)=-1$,
\inserted{whence:}
\[
u(\ell_\infty)=-u(\ell_1\ell_0)^{-1}
=-u(\ell_0^{-1})u(\ell_1)^{-1}
\]
we also have $u(\ell_\infty)\in\Pi$, OK.

If we identify $u$ with an element of $\mathcal M_{0,3}^{\sim}$,
we see that
\deleted{\unsure{it is}} in fact
\deleted{an element of $\unsure{Z|\mathcal M}$.}

\mspage{580}

% Printed archive footer outside the manuscript frame (not manuscript):
% 580
% https://grothendieck.umontpellier.fr
% Université de Montpellier, UMR CNRS 5149, Case courrier 051, Place Eugène Bataillon
% 34095 Montpellier cedex 5 - France

\scanpage{582}

% source-position: upper portion, continuation from preceding page.
\[
u\in\mathcal M_{0,3}^{\sim}[0]\simeq\mathcal M^![0]\subset\mathcal M
\]
% source-position: large slanted upper-left margin, correction and later-page reference.
\margin{Careful
we had not previously introduced
$\unsure{\mathcal M_{0,3}^{\sim}|\mathcal M_{0,3}}[0]$, but only
$\mathcal M_{0,3}^{!\sim}[0]\subset\mathcal M_{0,3}^{!\sim}$.
Now it \unsure{is not}
true that we shall have here $u\in\mathcal M_{0,3}^{!\sim}$,
for example \unsure{we may have} $u=\varepsilon_0$.
The image of $u$ in $\mathcal G_3$ is in
\deleted{\unsure{\operatorname{sg}}}
$\{1,\dot\sigma_0\}$, thus by correcting
\inserted{by $\unsure{\varepsilon_0}$},
we shall have an element of $\mathcal M_{0,3}^{!\sim}[0]$.
\deleted{We must}
Cf. more detailed correction on page $609$ ff.
$Z_{\rho,\sigma}^!=\mathcal M^![0]=Z_{\rho,\sigma}\cap\mathcal M^!$.}

% source-position: upper portion, display (38) and its cancelled side conditions.
\deleted{\unsure{we have}} thus $(20)$, more precisely
\[
(38)\qquad
\left\{
\begin{aligned}
&u(\rho)=\operatorname{int}(\alpha)\rho\\
&u(\sigma)=\varepsilon_2\operatorname{int}(\beta)\sigma\\
&u(\varepsilon_0)=\varepsilon_0^\mu
\end{aligned}
\right.
\qquad
\begin{aligned}
&\unsure{\text{with}}\ \deleted{\unsure{\alpha}}
\quad\alpha\in\widehat{\Pi}_{0\tau},\ \beta\in\widehat{\Pi}_0\\
&\deleted{\unsure{\varepsilon_2}}\quad
\left\{\begin{aligned}
&\mu\in\widehat{\mathbb Z}^{*}\\
&\deleted{\varepsilon_2=\varepsilon_2(\mu),\
\varepsilon_\tau(\alpha)=\varepsilon_3(\mu)}
\end{aligned}\right.
\end{aligned}
\]
% source-position: upper-middle portion, uniqueness paragraph after (38).
\unsure{in} their turn \unsure{imply|implying} $(37)$. In fact
the formulas \inserted{$(38)$} for $u(\rho),u(\sigma),u(\varepsilon_0)$
uniquely determine
$\mu\in\widehat{\mathbb Z}^{*}$, $\varepsilon_2\in\{\pm1\}$,
$\alpha\in\widehat{\Pi}_{0,\tau}$, $\beta\in\widehat{\Pi}_0$,
and we shall then necessarily have
% source-position: middle portion, display (39), exact sequence (40).
\[
(39)\qquad
\varepsilon_2=\varepsilon_2(\mu),\qquad
\varepsilon_\tau(\alpha)=\varepsilon_3(\mu)
\]
where $\varepsilon_\tau$ is the homomorphism that occurs in the
\inserted{exact} sequence
\[
(40)\qquad
\deleted{\unsure{\varepsilon_0}}
\diagram{
\begin{tikzcd}
1\arrow[r]&\widehat{\Pi}_0\arrow[r]&
\widehat{\Pi}_{0,\tau}\arrow[r,"\varepsilon_\tau"]&
\{\pm1\}\arrow[r]&1
\end{tikzcd}
}.
\]

% source-position: lower-middle portion, proposition and component identifications.
Thus we have proved

\uline{Proposition}\quad
In the universal case $(\mathcal G=\mathcal G^{+\wedge})$
we have
\[
\diagram{
\begin{tikzcd}
Z_{\rho,\sigma}^!&
\mathcal M^![0]\stackrel{\mathrm{dfn}}=
\mathcal M_{0,3}^{!\sim}[0]\arrow[l,"\sim"']
\end{tikzcd}
},
\]
the components $u,\mu,\varepsilon',\varepsilon$ of an element of
$Z_{\rho,\sigma}^!$ correspond resp.
(in terms of $u_0\in\mathcal M[0]$):
$u_0$ identified with an element of
$\operatorname{Aut}(\mathcal G^{+\wedge})$,
$\mu=\chi(u_0)$ (the multiplier of $u_0$),
$\varepsilon=\varepsilon_3(\mu)$,
$\varepsilon'=\varepsilon_2(\mu)$.

% source-position: lower portion, corollary and chain of quotients.
\uline{Corollary}\quad We have
\[
G_{\rho,\sigma}\simeq\mathcal M=\mathcal N_{1,1}^{\sim},
\]
\[
\gamma_{\rho,\sigma}\simeq
\mathcal M/\mathcal G^{+\wedge}
\simeq\mathcal N_{1,1}^{\sim}/\mathcal G^{+\wedge}
\simeq\mathcal M_{0,3}^{\sim}/\mathcal G_{0,3}^{+\wedge}
\simeq\Gamma_{\mathbb Q}^{\sim}.
\]

\mspage{581}

% Printed archive footer outside the manuscript frame (not manuscript):
% 581
% https://grothendieck.umontpellier.fr
% Université de Montpellier, UMR CNRS 5149, Case courrier 051, Place Eugène Bataillon
% 34095 Montpellier cedex 5 - France

\scanpage{583}
\mspage{595}
% source-position: upper heading, Remarque and (41)-(42).

\uline{Remark} It is tempting to write the \inserted{first two}
relations (38) in the form
\[
(41)\qquad
\left\{\begin{array}{ll}
(\alpha^{-1}u)(\rho)=\rho
  & \text{i.e.}\quad \alpha^{-1}u\in Z_\rho^{\unsure{\sharp}}\\
(\beta^{-1}u)(\sigma)=\varepsilon_2\sigma
  & \unsure{\text{whence}|\text{i.e.}}\quad \beta^{-1}u\in N'_\sigma\\
\deleted{\unsure{u(\varepsilon_0)}} &
\end{array}\right.
\]
and then to introduce
\[
(42)\qquad
\left\{\begin{array}{ll}
a=\alpha^{-1}u\in Z_\rho^{\unsure{\sharp}},
  & \text{thus }u=\alpha a,\quad \alpha=ua^{-1}\\
b=\beta^{-1}u\in N'_\sigma,
  & \text{thus }u=\beta b,\quad \beta=ub^{-1},
\end{array}\right.
\]
% source-position: middle paragraph and quotient diagram preceding (43).
and the developments of \S\,48 (notably
\uline{\(\overline{\mathrm{XI}}\)}) may lead one to expect that the
quantities \(a,b\), for \(u\in\mathcal{M}[0]\) \unsure{variable},
depend \uline{multiplicatively} on \(u\). We are going to see however
that this is not so. Let indeed
\diagram{\begin{tikzcd}[column sep=large]
\delta:\ \deleted{\unsure{\mathcal{M}}}\ G_{\rho,\sigma}
  \arrow[r,"\sim"] & \mathcal{M}
  \arrow[r] & \gamma_{\rho,\sigma}\simeq\Gamma_{\mathbb Q}
\end{tikzcd}}
be the canonical homomorphism, so that we will have
\[
(43)\qquad
\delta(u)=\delta(b)=\tau_\gamma^{\unsure{\nu_3}}\delta(a),
\]
putting
\[
\tau_\gamma=\delta(\tau),\qquad
\varepsilon_3=(-1)^{\nu_3}\quad
(\nu_3\in\mathbb Z/2\mathbb Z)
\]
% source-position: lower half, (43 bis), (43 ter), and (44).
If \(u'\) corresponds to \(b',a',\) \inserted{\(\nu'_3\)}, we will
likewise have
\[
(43\,\mathrm{bis})\qquad
\delta(u')=\delta(b')=\tau_\gamma^{\nu'_3}\delta(a')
\]
and \unsure{if} \(u''=u'u\) corresponds to \(a'',b'',\nu''_3\),
\[
(43\,\mathrm{ter})\qquad
\delta(u'u)=\delta(b'')=\tau_\gamma^{\nu''_3}\delta(a''),
\]
In fact, we will have
\[
\begin{array}{rl}
(44)\qquad & b''=b'b,\quad \nu''_2=\nu'_2+\nu_2\quad
 \unsure{\text{and.}|\text{and}}\quad \nu''_3=\nu'_3+\nu_3\\
& (\text{i.e. }\varepsilon''_2=\varepsilon'_2\varepsilon_2
\quad\text{and}\quad\varepsilon''_3=\varepsilon'_3\varepsilon_3)
\end{array}
\]
\mspage{582}

% Printed archive footer outside the manuscript frame (not manuscript):
% 582
% https://grothendieck.umontpellier.fr
% Université de Montpellier, UMR CNRS 5149, Case courrier 051, Place Eugène Bataillon
% 34095 Montpellier cedex 5 - France

\scanpage{584}
% source-position: upper calculations leading to (45).

We have seen
\[
\begin{array}{l}
\deleted{\unsure{\varepsilon_2|\text{and}}}u(\sigma)
 =\varepsilon_2\operatorname{int}(\beta)\sigma\\
u'(\sigma)=\varepsilon'_2\operatorname{int}(\beta')\sigma
\end{array}
\]
whence
\[
\begin{array}{l}
(u'u)(\sigma)=u'(u(\sigma))
 =\varepsilon_2\operatorname{int}(u'(\beta))u'(\sigma)\\
\hfill{}=\varepsilon_2\varepsilon'_2
 \operatorname{int}(\deleted{\unsure{\beta'}}u'(\beta)\unsure{\beta')\sigma}
\end{array}
\]
with \(u'(\beta)\beta'\in\widehat{\Pi}_0\), which by uniqueness proves
\[
(45)\qquad
\beta''=u'(\beta)\beta',\qquad
\varepsilon''_2=\varepsilon'_2\varepsilon_2
\]
% source-position: middle, b-coordinate multiplication and underbraces.
whence we indeed conclude
\[
\begin{array}{ccc}
b''&=&\underbrace{b'}_{\beta'^{-1}u'}\ \underbrace{b}_{\beta^{-1}u}\\
\Vert&&\\
\beta''^{-1}u''=\beta''^{-1}(u'u)&&
\end{array}
\]
\[
\begin{array}{lll}
\text{i.e.}&\beta''^{-1}u'u=\beta'^{-1}u'\beta^{-1}u
 &\text{i.e.}\quad\beta''^{-1}=\beta'^{-1}u'(\beta^{-1})\\
\text{i.e.}&\beta''=u'(\beta)\beta'
 &\text{which is none other than (45).}
\end{array}
\]

% source-position: lower middle, epsilon_3 identity and (46).
The relation \(\varepsilon''_3=\varepsilon'_3\varepsilon_3\) is also
evident, \unsure{for example} from the relation (which contains \(\varepsilon_3\),
and likewise \(\varepsilon'_3,\varepsilon''_3\)
\[
(46)\qquad
u(\rho)\deleted{\text{ is}}\ \mathcal{G}^{+\wedge}\text{-conjugate to }\rho^{\varepsilon_3}.
\]
\deleted{\unsure{when we}}

% source-position: lower third, hypothetical a-coordinate law, (*) and final displays.
It remains therefore to examine the hypothetical relation
\[
a''=a'a
\]
which \deleted{\unsure{is written|would be}} implies
\[
(*)\qquad\delta(a'')=\delta(a')\delta(a)
\]
but comparing (43) (43 bis) (43 ter) we find the relation
\[
\deleted{\unsure{\delta}}\delta(u')\delta(u)
 =\tau_\gamma^{\nu'_3}\delta(a')\tau_\gamma^{\nu_3}\delta(a)
\]
\[
\delta(u'\unsure{u''|u})=\tau_\gamma^{\nu''_3}\delta(a'')
\]
\hfill i.e.
\mspage{583}

% Printed archive footer outside the manuscript frame (not manuscript):
% 583
% https://grothendieck.umontpellier.fr
% Université de Montpellier, UMR CNRS 5149, Case courrier 051, Place Eugène Bataillon
% 34095 Montpellier cedex 5 - France

\scanpage{585}
\deleted{\(\unsure{\tau_\gamma^{\nu''_3}}\)}
\mspage{596}
% source-position: upper cancelled sentence, identities and (47).

\deleted{so that the relation (*) is equivalent to}
\[
\tau_\gamma^{\nu'_3+\nu_3}\delta(a'')
 =\tau_\gamma^{\nu'_3}\delta(a')\tau_\gamma^{\nu_3}\delta(a)
\]
or again
\[
\deleted{\tau_\gamma^{\unsure{\nu_3}}}\delta(a'')
 =\deleted{\unsure{\tau_\gamma}}\inserted{\tau_\gamma^{\nu_3}}
 \delta(a')\tau_\gamma^{\nu_3}\delta(a)
 \deleted{\unsure{\tau_\gamma^{\nu_3}}}
\]
i.e.
\[
(47)\qquad
\delta(a'')=\deleted{\unsure{\operatorname{int}}}
 \tau_\gamma^{\nu_3}(\delta(a'))\cdot\delta(a).
\]
% source-position: upper middle, (48) and its left marginal note.
Thus the relation (*) is equivalent to
\[
(48)\qquad
\begin{array}{lll}
\delta(a')=\tau_\gamma^{\nu_3}(\delta(a'))
 &\text{i.e.}&[\delta(a'),\tau_\gamma^{\nu_3}]=1\quad
 \deleted{\unsure{\text{or}}}\\
 &\text{i.e.}&[\delta(u'),\tau_\gamma^{\nu_3}]=1\quad
 \deleted{\unsure{\text{(for}}}
\end{array}
\]
\margin{since \(\delta(u)=\tau_\gamma^{\nu_3}\delta(a)\)}

% source-position: middle, commutation and heuristic comparison with Gamma_Q.
This shows that if \(\nu_3=1\) i.e. \(\varepsilon_3=-1\), then
the relation (*) is equivalent to the \uline{commutation}
\deleted{\unsure{as}}
\uline{of} \(\delta(u')\) \deleted{\(\unsure{\delta(a')}\)}
\uline{with} \(\tau_\gamma\),
\deleted{\unsure{commutation i.e. to}}.
It is a draconian condition on an element of \(\Gamma_{\mathbb Q}^{\sim}\)
\unsure{that|which}, for an element of \(\Gamma_{\mathbb Q}\)
itself \unsure{in its version}, implies that
\(u'\in\{1,\tau_\gamma\}\) --- and by heuristic arguments, we have seen
that it should even be so for \unsure{every element}
of \(\Gamma_{\mathbb Q}^{\sim}\). It appears that the success of the presentation
of \S\,48 \uline{\(\overline{\mathrm{XI}}\)} is due to the fact that
\unsure{in} \unsure{the situations|the situation}
\unsure{considered|considered} in \unsure{the|the} cases,
the group \(\gamma_{\rho,\sigma}\simeq\widehat{\mathbb Z}^{*}\)
was \uline{commutative}, \inserted{or \(\tau_\gamma\) \uline{central}.}
% source-position: lower third, inclusion and restricted multiplicativity; cancelled left margin.
Outside this case, the quantity
\diagram{\begin{tikzcd}
a=\alpha^{-1}u\in Z_\rho^{\unsure{\sharp}}
 \arrow[r,hook] & G_{\rho,\sigma}
\end{tikzcd}}
depends multiplicatively on the element
\(u\in\deleted{\unsure{Z_{\rho,\sigma}}}\mathcal M^{!}[0]\),
only on condition of restricting oneself to \(u\in\mathcal M^{!}[0]'\).

\margin{\deleted{It even suffices that \(\tau_\gamma\)
be \unsure{centralized|central} by}
\diagram{\begin{tikzcd}[nodes in empty cells]
\deleted{\gamma^{!}_{\rho,\sigma}=\operatorname{Ker}(\gamma_{\rho,\sigma}}
 \arrow[r,"\deleted{\unsure{\varepsilon_3}}"'] & {}
\end{tikzcd}}}
\mspage{584}

% Printed archive footer outside the manuscript frame (not manuscript):
% 584
% https://grothendieck.umontpellier.fr
% Université de Montpellier, UMR CNRS 5149, Case courrier 051, Place Eugène Bataillon
% 34095 Montpellier cedex 5 - France

\scanpage{586}
% source-position: upper text and homomorphism diagram.

We are going to try to salvage things, by finding a canonical
\uline{homomorphism}
\diagram{\begin{tikzcd}
\mathcal M^{!}[0] \arrow[r] \arrow[d,equal]
 & \underbrace{Z_{\rho,\sigma}^{!}
 \times_{\gamma_{\rho,\sigma}}\deleted{\unsure{Z}}N_\rho^{*}
 \times_{\gamma_{\rho,\sigma}}N_\sigma^{*}}
 _{S_0\Gamma_{\rho,\sigma}^{!}\underset{\text{index }2}{\subset}
 S_0\Gamma_{\rho,\sigma}} \\
Z_{\rho,\sigma}^{!} &
\end{tikzcd}}

\margin{\uline{N.B.} \unsure{we have}
\diagram{\begin{tikzcd}
S_0\Gamma_{\rho,\sigma}^{!}/L_0^{\Pi}\arrow[r,"\sim"]
 & S_0\Gamma_{\rho,\sigma}/L_0^{\mathcal G}
\end{tikzcd}}
\uuline{then}
\diagram{\begin{tikzcd}
\gamma_{\rho,\sigma}\arrow[r,"\sim"] & \Gamma_{\mathbb Q}^{\sim}.
\end{tikzcd}}}

% source-position: upper quarter, annotated triples (49).
Let us note that the second \unsure{side|set} is formed of the triples
\diagram{\begin{tikzcd}[column sep=small,row sep=small]
(49)\quad\bigl\{\bigl(u,\arrow[d,phantom,"\in"{rotate=90}]
 & a,\arrow[d,phantom,"\in"{rotate=90}]
 & \deleted{\unsure{b}}b\bigr)\mid\arrow[d,phantom,"\in"{rotate=90}]
 & u=fb=f'a,\ \text{with }f,f'\in\mathcal G^{+\wedge}
 \text{ suitable}\bigr\} \\
\deleted{\unsure{Z}}\unsure{\mathcal M}[0]
 & N'_\rho & N'_\sigma &
\end{tikzcd}}
\[
\text{which implies}\quad
\begin{array}{l}
\varepsilon_2(u)=\varepsilon_\sigma(b)\\
\varepsilon_3(u)=\varepsilon_\rho(a)
\end{array}
\]
and for such a triple, we will have
% source-position: upper middle, calculations (49 bis) and definition.
\[
(49\ \mathrm{bis})\quad
\left\{\begin{array}{l}
u(\deleted{\unsure{\rho}}\inserted{\sigma})
 =\operatorname{int}(\unsure{fb})
 \deleted{\unsure{\varepsilon_2(b)|\operatorname{int}(b)}}(\sigma)
 =\operatorname{int}(f)\underbrace{\operatorname{int}(b)(\sigma)}
 _{\varepsilon_2(b)(\sigma)}
 =\deleted{\operatorname{int}(f)\unsure{\varepsilon_2}}
 \inserted{\varepsilon_2(u)\operatorname{int}(f)(\sigma)}\\
u(\rho)=\operatorname{int}(f')(\rho^{\varepsilon_3(u)})
 =\operatorname{int}(f')\varepsilon_3(\tau')(\rho)
 =\operatorname{int}(f'\varepsilon_3(\tau'))(\rho)
\end{array}\right.
\]
where we put
\[
\varepsilon_3(\tau')\overset{\mathrm{dfn}}{=}
\left\{\begin{array}{ll}
1&\text{if }\varepsilon_3=+1\\
\unsure{\tau'}&\text{if }\varepsilon_3=-1
\end{array}\right.
\qquad
\begin{array}{l}
\text{i.e. }\varepsilon_3(\tau')=\sigma^{\nu_3}\\
\text{if }\varepsilon_3=(-1)^{\nu_3}
\end{array}
\]

If \(u\) corresponds to \(\alpha,\beta\) (cf. (38)), the
% source-position: middle, equivalent conjugation formulas after (49 bis).
formulas (49 bis) are written \unsure{also|thus} respectively
\[
\left\{\begin{array}{lll}
\operatorname{int}(f)(\sigma)=\operatorname{int}(\beta)(\sigma)
 &\text{i.e.}& f=\beta\sigma^i\quad(i\in\mathbb Z/4\mathbb Z)\\
\deleted{\operatorname{int}(f')\unsure{\varepsilon_3}(\rho)=}&&\\
\operatorname{int}(f'\varepsilon_3(\tau'))(\rho)
 =\operatorname{int}(\alpha)(\rho)
 &\text{i.e.}&\underbrace{f'\varepsilon_3(\tau')=\alpha\rho^j}
 _{j\in\mathbb Z/6\mathbb Z}
\end{array}\right.
\]

Thus we find

% source-position: lower third, Proposition and (50).
\uline{Proposition}
For a fixed \(u\in Z_{\rho,\sigma}^{!}=\mathcal M[0]
\simeq\unsure{\widetilde{\mathcal M}_{0,3}}[0]\),
the \((u,a,b)\in S_0\Gamma_{\rho,\sigma}\) over \(u\)
are exactly the elements \deleted{of the form}\inserted{of}
\[
(50)\quad
\left\{(u,a,b)\left|\begin{array}{lll}
a=f'^{-1}u&\inserted{\text{with}}\quad
 f'=\alpha\rho^j\varepsilon_3(\tau')&(j\in\mathbb Z/6\mathbb Z)\\
b=f^{-1}u&\text{with }f=\beta\sigma^i&(i\in\mathbb Z/4\mathbb Z)
\end{array}\right.\right\}
\]
\margin{where \(\alpha,\beta,\varepsilon_2,\varepsilon_3\) are
defined for \(u\) \unsure{via} (38)}

\deleted{We will have}\inserted{Let us consider} then the element
% source-position: final paragraph following (50), ending at the bottom of the scan.
\inserted{\((u,a,b)\)} over \(u\)
having the simplest expression in terms of \(\alpha,\beta\),
by taking \(i=\unsure{0|j},\ j=0\), thus corresponding to the choices:
\mspage{585}

% Printed archive footer outside the manuscript frame (not manuscript):
% 585
% https://grothendieck.umontpellier.fr
% Université de Montpellier, UMR CNRS 5149, Case courrier 051, Place Eugène Bataillon
% 34095 Montpellier cedex 5 - France
